\subsection{Stage~B construction comparison}
\label{sec:results-stageB}

This section compares the per-threshold and max-type Stage~B rules on size,
power, threshold placement, cost, and the benchmark rankings. Both rules
control the false-split rate. At matched realized size the max-type rule has
slightly lower power, and it places the threshold slightly less precisely. It
is cheaper once a node has many candidates.

The complete-node false-split rate is below 0.05 for both rules in all 36
design cells (\cref{tab:stageB-calibration}, \cref{app:stageB-tables}). These
rates include the Stage~A test over five predictors, so they bound the
threshold test from below rather than measure it.

The Stage-B-only study measures the threshold test itself
(\cref{tab:stageB-only}). At nominal level 0.05 the per-threshold rule has size
0.000 to 0.027, with means 0.018, 0.011, and 0.006 at 16, 64, and 256 bins.
The max-type rule has size 0.023 to 0.044, with means 0.033, 0.035, and 0.035.
Every 95\% Clopper-Pearson upper limit is below 0.056. Both rules are valid at
a fixed node for a predictor fixed without the response. The reverse
cardinality bias of \cref{sec:rank-construction} appears as a threefold fall in
per-threshold size across bin settings, while the max-type size does not move.

\begin{table}[!htbp]
  \centering
  \footnotesize
  \setlength{\tabcolsep}{4pt}
  \begin{tabular}{@{}lrrrrrrrrr@{}}
    \toprule
     & & & \multicolumn{4}{c}{Size at nominal 0.05} & & \multicolumn{2}{c}{Size-matched power} \\
     & & & \multicolumn{2}{c}{Per-threshold} & \multicolumn{2}{c}{Max-type} & Match & & \\
    \tablehead{Task} & $n$ & Bins & Adaptive & None & Adaptive & None & size & Per-threshold & Max-type \\
    \midrule
    Classification & 100 & 16 & 0.013 & 0.013 & 0.032 & 0.033 & 0.050 & 0.390 & 0.379 \\
     & 100 & 64 & 0.007 & 0.007 & 0.036 & 0.039 & 0.044 & 0.385 & 0.354 \\
     & 100 & 256 & 0.007 & 0.007 & 0.036 & 0.042 & 0.034 & 0.335 & 0.318 \\
     & 200 & 16 & 0.010 & 0.011 & 0.037 & 0.039 & 0.050 & 0.643 & 0.618 \\
     & 200 & 64 & 0.008 & 0.009 & 0.035 & 0.039 & 0.050 & 0.630 & 0.615 \\
     & 200 & 256 & 0.003 & 0.003 & 0.039 & 0.039 & 0.031 & 0.562 & 0.538 \\
     & 500 & 16 & 0.015 & 0.015 & 0.032 & 0.037 & 0.050 & 0.958 & 0.953 \\
     & 500 & 64 & 0.007 & 0.007 & 0.029 & 0.035 & 0.050 & 0.961 & 0.964 \\
     & 500 & 256 & 0.000 & 0.001 & 0.023 & 0.032 & 0.034 & 0.958 & 0.952 \\
    Regression & 100 & 16 & 0.019 & 0.019 & 0.023 & 0.031 & 0.050 & 0.292 & 0.281 \\
     & 100 & 64 & 0.013 & 0.013 & 0.029 & 0.035 & 0.050 & 0.305 & 0.285 \\
     & 100 & 256 & 0.011 & 0.010 & 0.031 & 0.037 & 0.050 & 0.317 & 0.299 \\
     & 200 & 16 & 0.027 & 0.026 & 0.033 & 0.038 & 0.050 & 0.591 & 0.581 \\
     & 200 & 64 & 0.017 & 0.017 & 0.033 & 0.043 & 0.050 & 0.619 & 0.598 \\
     & 200 & 256 & 0.008 & 0.009 & 0.041 & 0.044 & 0.050 & 0.632 & 0.579 \\
     & 500 & 16 & 0.026 & 0.026 & 0.031 & 0.032 & 0.050 & 0.947 & 0.942 \\
     & 500 & 64 & 0.015 & 0.016 & 0.029 & 0.033 & 0.050 & 0.966 & 0.952 \\
     & 500 & 256 & 0.005 & 0.006 & 0.026 & 0.034 & 0.046 & 0.963 & 0.951 \\
    \bottomrule
  \end{tabular}
  \caption{Stage-B-only study: the threshold test in isolation, with one
  Gaussian predictor fixed without the response and the Stage~A test disabled.
  Size is the false-split rate of a depth-one tree at nominal level 0.05 over
  1,500 null replicates per cell (three seeds of 500) under adaptive stopping
  and no stopping (None); every 95\% Clopper-Pearson upper limit is below
  0.056. Size-matched power, shown for adaptive stopping at the middle effect
  ($\delta = 0.10$ in classification, $0.4$ in regression), is interpolated
  against realized size over the nominal grid from 0.02 to 0.50 (600 replicates
  per cell and level) to the match size: 0.05, or the largest realized size
  both rules reach where the per-threshold rule does not reach 0.05.}
  \label{tab:stageB-only}
\end{table}

Matched at the same realized size, the max-type minus per-threshold difference
in power ranges from $-0.054$ to $+0.034$ over the 108 cell and effect
combinations. The max-type rule is lower in 67 and higher in 25. The mean
difference is $-0.006$, with a 95\% bootstrap interval of $-0.011$ to $-0.002$
that resamples the 18 cells, and $-0.005$ ($-0.009$ to $-0.002$) over the 13
cells matched at size 0.05. It has no trend in the bin setting. In the other
five cells the per-threshold rule does not reach size 0.05 even at nominal
0.50, and the match is at the largest size both rules reach.

Two reruns bound the Monte Carlo error. With larger replicate counts at a
fourth seed, the mean difference is $-0.006$ ($-0.008$ to $-0.003$); a
parametric bootstrap of the whole matching procedure puts a standard error of
about 0.017 on each difference, so the smallest detectable difference is about
0.03. With the planted step moved to the 85th percentile, the mean difference
is $-0.004$ ($-0.008$ to $0.000$), with a standard error of about 0.038 and a
detectable difference of about 0.08. The Stage-B-only sizes stay in the ranges
above.

At a common nominal level of 0.05, the max-type rule splits more often, by 0.03
on average and up to 0.12 in the complete-node design (\cref{tab:stageB-power})
and by 0.10 in the Stage-B-only study. The gap grows with the candidate count,
0.015 at 16 bins, 0.029 at 64, and 0.049 at 256, because the per-threshold
power falls as the threshold correction tightens. The gap therefore reflects
the lower realized size of the per-threshold rule, which the matched
comparison above removes.

The max-type rule places the threshold less precisely. On each rule's own
detections, the max-type threshold sits a median 0.19 standard deviations from the planted cut
against 0.15 for the per-threshold rule, and 87\% of its detections fall on
the planted predictor against 90\%.

\Cref{tab:stageB-scaling} reports the cost of the two rules. With no stopping and
256 bins, one classification tree under the per-threshold rule takes 3.5, 14,
and 59 seconds at $n=1{,}000$, $4{,}000$, and $16{,}000$ (0.8, 8.8, and 54 in
regression); the max-type rule takes 0.3, 1.4, and 6.3 seconds. The ratio grows
with the bin setting as the $K^2/\alpha$ against $BK$ accounting predicts. At
16 bins the max-type rule is 1.1 to 1.3 times slower, because its 999
permutations exceed the 340 the per-threshold rule needs for 17 candidates. At
64 bins it is 1.5 to 2.1 times faster, and at 256 bins 2.8 to 13 times.
Adaptive stopping shortens the per-threshold test by 1.2 to 1.6, 2.2 to 3.2,
and 3.8 to 7.8 times at the three bin settings, so with adaptive stopping the
max-type advantage is 1.0 to 1.3, 1.4 to 2.1, and 2.2 to 7.7. Held at 999
permutations per candidate instead, the per-threshold rule never splits at 64
or 256 bins, because $\alpha/K$ falls below the resolution $1/(B+1)$. Fitted
depths barely differ: 5.0 to 10.7 levels for the max-type rule against 4.0 to
10.0.

On the six Stage~B real datasets (\cref{tab:stageB-real}), the max-type rule
scores higher in 8 of the 12 dataset and stopping-mode cells and equal in one.
It is higher by 0.016 and 0.026 balanced accuracy on vowel, 0.049 and 0.026
$R^2$ on facebook, 0.060 and 0.033 on imports-85, and 0.008 and 0.005 on
residential. It is lower by 0.018 and 0.020 on page-blocks, and on spam lower
by 0.006 with adaptive stopping and equal with no stopping. Its fitting time is
1.2 to 2.4 times shorter with adaptive stopping and 1.0 to 8.3 times shorter
with no stopping on the four datasets whose nodes admit many candidates. It is
16\% and 10\% slower on residential. Fitted depths agree within about one
level (5.4 to 11.4 for the per-threshold rule, 4.8 to 12.4 for the max-type
rule). The Spearman agreement of the two rules' feature importances is 0.36 to
0.88, so the max-type rule produces a different ranking rather than the same
ranking at lower cost.

In the head-to-head forest timing, the max-type rule is 10 to 27 times faster
than the per-threshold rule on the synthetic cells at $n\le 8{,}000$ and 21 to
36 times faster at $n=32{,}000$ (13 to 36 times on the cells that
\citet{MilletichDownesHirst2026citrees} tabulate). The runtime panel of
\cref{tab:runtime-ablations} shows smaller ratios, 1.6 to 7 times, because its
nodes admit fewer candidates. On the ten real datasets the max-type rule
reduces madelon from 21 to 2.6 seconds, gamma from 44 to 21, and isolet from 18
to 9.1, and changes the other seven by at most 1.6 seconds.

The max-type extension shows what the rule does to the rankings
(\cref{tab:maxt-benchmark-rankings}). Of its 1,880 possible ranking runs, 88
were not run because their per-threshold counterpart had not completed, and
1,785 of the remaining 1,792 completed. The seven that did not are CIF with the
RDC selector and honesty off on isolet, gisette, and letter, the same cells
excluded from the per-threshold benchmark (\cref{app:methods}). Ranked on the
all-method panels, the max-type forest moves from 3.5 to 5 of 17 in
classification and from 2 to 4 of 18 in regression, and the max-type tree
keeps its position in both tasks. The per-threshold rule is the higher of the
two on 59 to 63\% of datasets. The case for the max-type rule therefore rests
on cost and calibration rather than on the rankings.

\begin{table}[!htbp]
  \centering
  \footnotesize
  \setlength{\tabcolsep}{5pt}
  \begin{tabular}{@{}llrrrrrr@{}}
    \toprule
     & & \multicolumn{2}{c}{Per-threshold} & \multicolumn{2}{c}{Max-type} & \multicolumn{2}{c}{Paired difference} \\
    \tablehead{Task} & Method & Mean rank & Position & Mean rank & Position & Score & Datasets \\
    \midrule
    Classification & CIF & 5.86 & 3.5 & 6.31 & 5 & $-0.0014$ & 22 \\
     & CIT & 13.14 & 16 & 13.00 & 16 & $-0.0022$ & 23 \\
    Regression & CIF & 6.25 & 2 & 7.50 & 4 & $-0.0001$ & 8 \\
     & CIT & 11.13 & 14 & 10.88 & 14 & $-0.0085$ & 8 \\
    \bottomrule
  \end{tabular}
  \caption{Benchmark rankings under the two Stage~B rules. Mean rank and
  position are among the 17 classification and 18 regression methods on the
  all-method panels of 21 and 8 datasets, with the max-type family standing in
  for its per-threshold counterpart; the per-threshold columns repeat
  \cref{tab:method-main-comparison}. The paired difference is the max-type
  score minus the per-threshold score (balanced accuracy in classification,
  $R^2$ in regression), averaged over the datasets, downstream learners, and
  standard $k$ that both rules completed; Datasets is the number of datasets in
  that average.}
  \label{tab:maxt-benchmark-rankings}
\end{table}

\FloatBarrier

\paragraph{Cardinality in the ranking.}
The weak-signal designs place 10 informative Gaussian features among
500-level, binary, and Gaussian noise columns
(\cref{tab:weak-signal-cardinality}). A ranker indifferent to cardinality would
give each of the 500-level and binary noise types 25/60 of the top-ten places
not taken by informative features. The designs can reveal two biases, and both
appear. The first is the preference of impurity importance for features with
many possible thresholds; the second is the reverse cardinality bias of the
per-threshold rule (\cref{sec:rank-construction}). At the weakest signal, random forest impurity importance fills
0.38 of its top ten with 500-level noise in both tasks, against a reference of
0.24, and none with binary noise. Per-threshold CIF is close to the reference
in classification (0.31 binary and 0.26 500-level against 0.29) and leans
toward binary noise in regression (0.40 binary and 0.24 500-level against
0.31). The max-type rule brings the binary share to 0.26 in classification and
0.30 in regression, against references of 0.28 and 0.27.

A ranking with less cardinality preference does not recover more informative
features in these designs. Random forest impurity
importance places the most informative features in its top ten at every signal
strength, and the downstream scores of all rankers are within 0.02 of each
other at every signal. The same forest ranked by out-of-bag permutation
importance keeps the preference, with a 500-level share above that of
impurity importance at the two weaker signals in both tasks. The benefit of a
ranking closer to cardinality-neutral appears on predictors whose cardinality
carries no signal, as in the clinical application of
\citet{MilletichDownesHirst2026citrees}, rather than in recovery of weak
marginal effects.

\begin{table}[!htbp]
  \centering
  \footnotesize
  \setlength{\tabcolsep}{5pt}
  \begin{tabular}{@{}lrlrrrrr@{}}
    \toprule
     & & & \multicolumn{3}{c}{Share of the top ten} & & \\
    \tablehead{Task} & Signal & Ranking & Informative & 500-level noise & Binary noise & Support & Score \\
    \midrule

    Classification & 0.05 & CIF & 0.31 & 0.26 & 0.31 & 37 & 0.52 \\
     & & CIF, max-type & 0.34 & 0.28 & 0.26 & 46 & 0.51 \\
     & & RF impurity & 0.42 & 0.38 & 0.00 & 70 & 0.52 \\
     & & RF permutation & 0.30 & 0.47 & 0.06 & 34 & 0.51 \\
     & 0.10 & CIF & 0.64 & 0.16 & 0.13 & 45 & 0.56 \\
     & & CIF, max-type & 0.66 & 0.20 & 0.06 & 53 & 0.55 \\
     & & RF impurity & 0.72 & 0.19 & 0.00 & 70 & 0.56 \\
     & & RF permutation & 0.49 & 0.36 & 0.02 & 36 & 0.56 \\
     & 0.20 & CIF & 0.95 & 0.04 & 0.01 & 59 & 0.68 \\
     & & CIF, max-type & 0.96 & 0.04 & 0.00 & 65 & 0.69 \\
     & & RF impurity & 0.99 & 0.01 & 0.00 & 70 & 0.69 \\
     & & RF permutation & 0.87 & 0.09 & 0.00 & 39 & 0.67 \\
    Regression & 0.05 & CIF & 0.26 & 0.24 & 0.40 & 17 & $-$0.11 \\
     & & CIF, max-type & 0.34 & 0.22 & 0.30 & 21 & $-$0.11 \\
     & & RF impurity & 0.43 & 0.38 & 0.00 & 70 & $-$0.11 \\
     & & RF permutation & 0.37 & 0.43 & 0.05 & 34 & $-$0.10 \\
     & 0.10 & CIF & 0.64 & 0.11 & 0.20 & 24 & $-$0.05 \\
     & & CIF, max-type & 0.65 & 0.13 & 0.15 & 28 & $-$0.04 \\
     & & RF impurity & 0.73 & 0.17 & 0.00 & 70 & $-$0.04 \\
     & & RF permutation & 0.57 & 0.27 & 0.01 & 37 & $-$0.06 \\
     & 0.20 & CIF & 0.95 & 0.03 & 0.02 & 36 & 0.28 \\
     & & CIF, max-type & 0.95 & 0.02 & 0.02 & 44 & 0.28 \\
     & & RF impurity & 0.99 & 0.01 & 0.00 & 70 & 0.29 \\
     & & RF permutation & 0.95 & 0.03 & 0.00 & 39 & 0.28 \\
        \bottomrule
  \end{tabular}
  \caption{Weak-signal designs with mixed-cardinality noise, $n = 1{,}000$.
  Ten informative standard normal features have marginal correlation with the
  response given by the Signal column; the noise is 25 columns of integer codes
  on 500 levels, 25 binary columns, and 10 Gaussian columns. Shares are the mean
  composition of each ranking's top ten over 25 fits (5 seeds by 5 folds;
  standard errors 0.01 to 0.03); the remainder is Gaussian noise. Support is
  the mean number of features with nonzero importance; Score is the held-out
  balanced accuracy or $R^2$ of the downstream learners on the top ten. CIF is
  the selected configuration with the MC selector in classification and the PC
  selector in regression; CIF, max-type adds the max-type Stage~B rule; RF
  permutation is the same random forest ranked by out-of-bag permutation
  importance.}
  \label{tab:weak-signal-cardinality}
\end{table}
